%
% lv0.tex -- Liouville-Form Basisform
%
% (c) 2021 Prof Dr Andreas Müller, OST Ostschweizer Fachhochschule
%
\documentclass[tikz]{standalone}
\usepackage{amsmath}
\usepackage{times}
\usepackage{txfonts}
\usepackage{pgfplots}
\usepackage{csvsimple}
\usetikzlibrary{arrows,intersections,math}
\begin{document}
\def\skala{1}
\begin{tikzpicture}[>=latex,thick,scale=\skala]

\draw (0,0.5) rectangle (1,3);
\node at (1,3) [below left] {$K\mathstrut$};

\node at (0.5,1) {$c_i\mathstrut$};

\draw (-0.05,0.45) rectangle (4,3.05);
\node at (4,3) [below left] {$F\mathstrut$};

\node at (2.5,1) {$v_0, v_i\mathstrut$};
\node at (2.5,2) {$f\mathstrut$};

\draw (-0.1,0.4) rectangle (7,3.1);
\node at (7,3) [below left] {$F(\vartheta)\mathstrut$};

\node at (5.5,2) {$g\mathstrut$};
\draw[double,line width=0.5pt, shorten >= 0.2cm,shorten <= 0.2cm] (5.5,2) -- (5.5,1);
\node at (4.5,2) [above] {$D\mathstrut$};
\node at (5.5,1) {$v_0+\sum_i c_i\log v_i\mathstrut$};

\draw[->,line width=0.5pt] (5.2,2) -- (2.8,2);
\draw[line width=0.5pt] (5.2,1.93) -- ++(0,0.14);

\end{tikzpicture}
\end{document}

